\setstretch{1.4}
\usepackage[usenames,dvipsnames]{color}
\usepackage[unicode=true, 
 bookmarks=true,bookmarksnumbered=false,bookmarksopen=true,bookmarksopenlevel=2,
 breaklinks=false,pdfborder={0 0 0},backref=false,colorlinks=false]
 {hyperref}
\hypersetup{pdftitle={Vergleich und Evaluation zwischen modernen und traditionellen Datenbankkonzepten 
unter den Gesichtspunkten Skalierung, Abfragem�glichkeit und Konsistenz.},
 pdfauthor={Nils M. Petersohn}}
 
\makeatletter

%%%%%%%%%%%%%%%%%%%%%%%%%%%%%% LyX specific LaTeX commands.
\providecommand{\LyX}{L\kern-.1667em\lower.25em\hbox{Y}\kern-.125emX\@}
%% Because html converters don't know tabularnewline
\providecommand{\tabularnewline}{\\}

%%%%%%%%%%%%%%%%%%%%%%%%%%%%%% Textclass specific LaTeX commands.
\newenvironment{lyxcode}
{\par\begin{list}{}{
\setlength{\rightmargin}{\leftmargin}
\setlength{\listparindent}{0pt}% needed for AMS classes
\raggedright
\setlength{\itemsep}{0pt}
\setlength{\parsep}{0pt}
\normalfont\ttfamily}%
 \item[]}
{\end{list}}

%%%%%%%%%%%%%%%%%%%%%%%%%%%%%% User specified LaTeX commands.
%% Flexibles Seitenlayout
\usepackage[automark]{scrpage2}

%% Mehrspaltenlayout ermöglichen
\usepackage{multicol}

%% Unterst\"utzung f\"ur Farben
\usepackage{color}

%% Schönere Tabellen
\usepackage{booktabs, longtable}

%% Schönerer Blocksatz
\usepackage{microtype}
%% Mehr Platz zwischen Überschrift und Tabelle
\newcommand{\@ldtable}{}
\let\@ldtable\table
\renewcommand{\table}{ %
    \setlength{\@tempdima}{\abovecaptionskip} %
    \setlength{\abovecaptionskip}{\belowcaptionskip} %
    \setlength{\belowcaptionskip}{\@tempdima} %
    \@ldtable %
}

%% Verschiedene Symbole und Zeichen wie (c), ¤
\usepackage{textcomp}

%% Fehlerkorrektur f\"ur Marginalien
\usepackage{fixltx2e }%,mparhack

%% Deutsche Kurzfassung und englisches Abstract auf eine Seite
\renewenvironment{abstract}{
    \@beginparpenalty\@lowpenalty
        \begin{center}
            \normalfont\sectfont\nobreak\abstractname
        \end{center}
    \@endparpenalty\@M
}{
    \par
}

%% Alle Seiten vor dem Inhaltsverzeichnis sind römisch nummeriert
\pagenumbering{roman}
\let\myTOC\tableofcontents
\renewcommand\tableofcontents{
    \begin{spacing}{1.1}
    \myTOC
    \end{spacing}
    \clearpage
    \pagenumbering{arabic}
}

%% Kopfzeile um Logo erg\"anzen
\clearscrheadfoot
\ohead{\\\headmark}
\ihead{\includegraphics[scale=0.25]{pics/FH-Logo-only-gray}}%\pagemark}
\ofoot[\pagemark]{\pagemark}

%% Randnotizen anpassen
\setlength{\marginparwidth}{22mm}
\let \oldmarginpar = \marginpar
\renewcommand{\marginpar}[1]{%
    \-\oldmarginpar[\raggedleft\footnotesize\sf #1]%
        {\raggedright\footnotesize\sf #1%
    }}

%% Zitate am Kapitelanfang
%\usepackage{epigraph}
%\setlength{\epigraphwidth}{9cm}

%% Code-Block-Formatierung
\lstdefinestyle{default}{ %
    backgroundcolor={\color[rgb]{0.95,0.95,0.95}}, %
    basicstyle={\small\ttfamily}, %
    breaklines=true, %
    frame=l, %
    language={[Sharp]C}, %
    lineskip=-0.1pt, %
    numbers=left, %
    rulecolor={\color[rgb]{0.5,0.5,0.5}}, %
    xleftmargin={0.75cm}, %
    xrightmargin={0cm} %
}
\lstdefinelanguage{JavaScript} {
	morekeywords={
		break,const,continue,delete,do,while,export,for,in,function,
		if,else,import,in,instanceOf,label,let,new,return,switch,this,
		throw,try,catch,typeof,var,void,with,yield
	},
	sensitive=false,
	morecomment=[l]{//},
	morecomment=[s]{/*}{*/},
	morestring=[b]",
	morestring=[d]'
}

\lstset{
	frame=tb,
	framesep=5pt,
	basicstyle=\footnotesize\ttfamily,
	showstringspaces=false,
	keywordstyle=\ttfamily\bfseries\color{CadetBlue},
	identifierstyle=\ttfamily,
	stringstyle=\ttfamily\color{OliveGreen},
	commentstyle=\color{GrayBlue},
	rulecolor=\color{Gray},
	xleftmargin=5pt,
	xrightmargin=5pt,
	aboveskip=\bigskipamount,
	belowskip=\bigskipamount
}
\makeatother

\setlength{\parindent}{5pt} 
\setlength{\parskip}{1ex}